\textbf{Phần 1. Trắc nghiệm lựa chọn}
\setcounterpageref{ex}{0}

\begin{ex} %[4H2V07]
Khi kim phút của đồng hồ quay được $\frac{1}{4}$ vòng (tương ứng với $15$ phút) thì kim phút đã quét được một góc có số đo là
\choice
{\True $90^\circ$}
{$60^\circ$}
{$180^\circ$}
{$45^\circ$}
\loigiai{
Một vòng quay trọn vẹn của kim đồng hồ tương ứng với $360^\circ$.\\
Khi kim phút quay được $\frac{1}{4}$ vòng thì góc quét được là:
$$360^\circ : 4 = 90^\circ$$
\textbf{Đáp án của câu hỏi là A}
}
\end{ex}

\begin{ex} %[4H2V07]
Lúc 8 giờ đúng, góc bé tạo bởi kim giờ và kim phút có số đo là
\choice
{$90^\circ$}
{$150^\circ$}
{\True $120^\circ$}
{$240^\circ$}
\loigiai{
Lúc 8 giờ đúng, kim phút chỉ số 12, kim giờ chỉ số 8.\\
Khoảng cách ngắn hơn giữa hai kim gồm $12 - 8 = 4$ khoảng giờ.\\
Mỗi khoảng giờ tương ứng với: $360^\circ : 12 = 30^\circ$.\\
Số đo góc bé tạo bởi hai kim là:
$$4 \times 30^\circ = 120^\circ$$
\textbf{Đáp án của câu hỏi là C}
}
\end{ex}

\begin{ex} %[4H2V07]
Cho góc $\widehat{xOy} = 75^\circ$ và góc $\widehat{yOz} = 45^\circ$ (hai góc kề nhau tạo thành góc $\widehat{xOz}$). Góc $\widehat{xOz}$ là
\choice
{Góc vuông}
{\True Góc tù}
{Góc nhọn}
{Góc bẹt}
\loigiai{
Số đo của góc $\widehat{xOz}$ là:
$$\widehat{xOz} = 75^\circ + 45^\circ = 120^\circ$$
Vì $90^\circ < 120^\circ < 180^\circ$ nên góc $\widehat{xOz}$ là góc tù.\\
\textbf{Đáp án của câu hỏi là B}
}
\end{ex}

\begin{ex} %[4H2V07]
Lúc 9 giờ 30 phút, góc bé tạo bởi kim giờ và kim phút của đồng hồ có số đo là
\choice
{$90^\circ$}
{$100^\circ$}
{$120^\circ$}
{\True $105^\circ$}
\loigiai{
Lúc 9 giờ 30 phút, kim phút chỉ đúng vào số 6.\\
Kim giờ nằm ở vị trí chính giữa số 9 và số 10 (tương ứng vị trí $9{,}5$).\\
Khoảng cách giữa hai kim là:
$$9{,}5 - 6 = 3{,}5\text{ (khoảng giờ)}$$
Số đo góc bé tạo bởi hai kim là:
$$3{,}5 \times 30^\circ = 105^\circ$$
\textbf{Đáp án của câu hỏi là D}
}
\end{ex}

\begin{ex} %[4H2V07]
Một chiếc quạt giấy khi xòe ra hết cỡ tạo thành một góc bẹt. Nếu người ta chỉ mở quạt bằng $\frac{2}{3}$ độ mở tối đa đó thì góc tạo bởi hai nan ngoài cùng của quạt có số đo là
\choice
{\True $120^\circ$}
{$90^\circ$}
{$135^\circ$}
{$150^\circ$}
\loigiai{
Góc bẹt có số đo bằng $180^\circ$.\\
Khi mở bằng $\frac{2}{3}$ độ mở tối đa, góc tạo bởi hai nan ngoài cùng là:
$$180^\circ \times 2 : 3 = 120^\circ$$
\textbf{Đáp án của câu hỏi là A}
}
\end{ex}

\begin{ex} %[4H2V07]
Cho hình chữ nhật $ABCD$. Kẻ đường chéo $AC$. Biết góc $\widehat{BAC} = 35^\circ$. Số đo của góc $\widehat{CAD}$ là
\choice
{$45^\circ$}
{$65^\circ$}
{\True $55^\circ$}
{$35^\circ$}
\loigiai{
Vì góc $\widehat{BAD}$ là góc vuông của hình chữ nhật nên $\widehat{BAD} = 90^\circ$.\\
Ta có:
$$\widehat{CAD} = 90^\circ - \widehat{BAC} = 90^\circ - 35^\circ = 55^\circ$$
\textbf{Đáp án của câu hỏi là C}
}
\end{ex}

\begin{ex} %[4H2V07]
Từ 1 giờ đúng đến 1 giờ 20 phút cùng ngày, kim phút của đồng hồ đã quét được một góc có số đo là
\choice
{$90^\circ$}
{\True $120^\circ$}
{$60^\circ$}
{$150^\circ$}
\loigiai{
Khoảng thời gian từ 1 giờ đúng đến 1 giờ 20 phút là $20$ phút.\\
Mỗi phút kim phút quay được:
$$360^\circ : 60 = 6^\circ$$
Trong 20 phút, góc kim phút đã quét được là:
$$20 \times 6^\circ = 120^\circ$$
\textbf{Đáp án của câu hỏi là B}
}
\end{ex}

\begin{ex} %[4H2V07]
Cho góc $\widehat{AOB} = 90^\circ$. Kẻ tia $OC$ nằm giữa hai tia $OA$ và $OB$ sao cho góc $\widehat{AOC}$ gấp đôi góc $\widehat{BOC}$. Số đo của góc $\widehat{AOC}$ là
\choice
{$30^\circ$}
{$45^\circ$}
{$50^\circ$}
{\True $60^\circ$}
\loigiai{
Ta có $\widehat{AOC} + \widehat{BOC} = 90^\circ$.\\
Coi góc $\widehat{BOC}$ là 1 phần thì góc $\widehat{AOC}$ là 2 phần bằng nhau như thế.\\
Tổng số phần bằng nhau là:
$$1 + 2 = 3\text{ (phần)}$$
Số đo của góc $\widehat{AOC}$ là:
$$90^\circ : 3 \times 2 = 60^\circ$$
\textbf{Đáp án của câu hỏi là D}
}
\end{ex}

\begin{ex} %[4H2V07]
Hình tam giác $ABC$ có góc đỉnh $A$ là góc vuông. Biết góc đỉnh $B$ lớn hơn góc đỉnh $C$ là $20^\circ$. Số đo của góc đỉnh $B$ là
\choice
{\True $55^\circ$}
{$35^\circ$}
{$60^\circ$}
{$50^\circ$}
\loigiai{
Trong tam giác vuông $ABC$, tổng hai góc nhọn đỉnh $B$ và đỉnh $C$ là $90^\circ$:\\
$$\widehat{B} + \widehat{C} = 90^\circ$$
Theo đề bài: $\widehat{B} - \widehat{C} = 20^\circ$.\\
Số đo của góc đỉnh $B$ là:
$$(90^\circ + 20^\circ) : 2 = 55^\circ$$
\textbf{Đáp án của câu hỏi là A}
}
\end{ex}

\begin{ex} %[4H2V07]
Lúc mấy giờ đúng thì góc bé tạo bởi hai kim giờ và kim phút có số đo bằng một góc bẹt?
\choice
{12 giờ đúng}
{3 giờ đúng}
{\True 6 giờ đúng}
{9 giờ đúng}
\loigiai{
Góc bẹt có số đo bằng $180^\circ$.\\
Khi hai kim giờ và kim phút tạo thành góc $180^\circ$ thì hai kim nằm trên cùng một đường thẳng và chỉ về hai hướng đối diện nhau. Đó là thời điểm 6 giờ đúng.\\
\textbf{Đáp án của câu hỏi là C}
}
\end{ex}

\begin{ex} %[4H2V07]
Một bạn học sinh đứng quay mặt về hướng Bắc, sau đó bạn quay sang bên phải một góc có số đo bằng $135^\circ$. Hỏi lúc này bạn đang quay mặt về hướng nào?
\choice
{Hướng Đông}
{\True Hướng Đông Nam}
{Hướng Nam}
{Hướng Đông Bắc}
\loigiai{
Từ hướng Bắc quay sang phải $90^\circ$ là hướng Đông.\\
Từ hướng Đông quay tiếp sang phải thêm $135^\circ - 90^\circ = 45^\circ$ sẽ tới hướng chính giữa của hướng Đông và hướng Nam, tức là hướng Đông Nam.\\
\textbf{Đáp án của câu hỏi là B}
}
\end{ex}

\begin{ex} %[4H2V07]
Cho 4 góc nhọn có số đo bằng nhau và tổng số đo của 4 góc đó đúng bằng một góc bẹt. Số đo của mỗi góc nhọn đó là
\choice
{$30^\circ$}
{$40^\circ$}
{$60^\circ$}
{\True $45^\circ$}
\loigiai{
Một góc bẹt có số đo bằng $180^\circ$.\\
Số đo của mỗi góc nhọn là:
$$180^\circ : 4 = 45^\circ$$
\textbf{Đáp án của câu hỏi là D}
}
\end{ex}

\textbf{Phần 2. Trắc nghiệm đúng sai}
\setcounterpageref{ex}{0}

\begin{ex} %[4H2V07]
Xét tính đúng sai của các khẳng định sau về số đo góc:
\choiceTF
{\True Góc có số đo $89^\circ$ là góc nhọn.}
{Góc có số đo $91^\circ$ là góc bẹt.}
{\True Hai góc vuông có tổng số đo bằng một góc bẹt.}
{Góc có số đo lớn hơn góc nhọn luôn là góc tù.}
\loigiai{
\begin{itemize}
    \item Vì $89^\circ < 90^\circ$ nên đây là góc nhọn. Khẳng định a Đúng.
    \item Vì $90^\circ < 91^\circ < 180^\circ$ nên góc $91^\circ$ là góc tù (không phải góc bẹt). Khẳng định b Sai.
    \item Tổng số đo hai góc vuông: $90^\circ + 90^\circ = 180^\circ$ (bằng một góc bẹt). Khẳng định c Đúng.
    \item Góc lớn hơn góc nhọn có thể là góc vuông ($90^\circ$) hoặc góc bẹt ($180^\circ$), không nhất thiết phải là góc tù. Khẳng định d Sai.
\end{itemize}
\begin{itemize}
    \item a) \textbf{Đúng}.
    \item b) \textbf{Sai}.
    \item c) \textbf{Đúng}.
    \item d) \textbf{Sai}.
\end{itemize}
}
\end{ex}

\begin{ex} %[4H2V07]
Xét các góc tạo bởi kim giờ và kim phút trên mặt đồng hồ:
\choiceTF
{\True Lúc 2 giờ đúng, góc bé tạo bởi hai kim có số đo là $60^\circ$.}
{\True Lúc 4 giờ đúng, góc bé tạo bởi hai kim có số đo là $120^\circ$.}
{Lúc 10 giờ đúng, góc bé tạo bởi hai kim có số đo là $100^\circ$.}
{\True Lúc 6 giờ đúng, hai kim tạo thành một góc bẹt có số đo $180^\circ$.}
\loigiai{
\begin{itemize}
    \item Lúc 2 giờ: $2 \times 30^\circ = 60^\circ$. Khẳng định a Đúng.
    \item Lúc 4 giờ: $4 \times 30^\circ = 120^\circ$. Khẳng định b Đúng.
    \item Lúc 10 giờ: khoảng cách ngắn nhất là $12 - 10 = 2$ khoảng giờ, góc là $2 \times 30^\circ = 60^\circ$ (không phải $100^\circ$). Khẳng định c Sai.
    \item Lúc 6 giờ: hai kim thẳng hàng tạo thành góc bẹt $180^\circ$. Khẳng định d Đúng.
\end{itemize}
\begin{itemize}
    \item a) \textbf{Đúng}.
    \item b) \textbf{Đúng}.
    \item c) \textbf{Sai}.
    \item d) \textbf{Đúng}.
\end{itemize}
}
\end{ex}

\textbf{Phần 3. Trả lời ngắn}
\setcounterpageref{ex}{0}

\begin{ex} %[4H2V07]
Lúc 5 giờ đúng, góc bé tạo bởi kim giờ và kim phút của đồng hồ có số đo bằng bao nhiêu độ? (Chỉ ghi số).
\shortans{150}
\loigiai{
\textbf{Đáp số:} $150$\\
Khoảng cách giữa hai kim gồm 5 khoảng giờ:
$$5 \times 30^\circ = 150^\circ$$
}
\end{ex}

\begin{ex} %[4H2V07]
Trong một giờ, kim phút của đồng hồ quay được một góc có số đo bằng bao nhiêu độ? (Chỉ ghi số).
\shortans{360}
\loigiai{
\textbf{Đáp số:} $360$\\
Trong một giờ, kim phút quay trọn vẹn đúng 1 vòng tròn:
$$360^\circ$$
}
\end{ex}

\begin{ex} %[4H2V07]
Một góc tù có số đo lớn hơn góc vuông $45^\circ$. Hỏi góc tù đó có số đo bằng bao nhiêu độ?
\shortans{135}
\loigiai{
\textbf{Đáp số:} $135$\\
Số đo của góc tù đó là:
$$90^\circ + 45^\circ = 135^\circ$$
}
\end{ex}

\begin{ex} %[4H2V07]
Lúc 2 giờ 30 phút, góc bé tạo bởi kim giờ và kim phút của đồng hồ có số đo bằng bao nhiêu độ?
\shortans{105}
\loigiai{
\textbf{Đáp số:} $105$\\
Lúc 2 giờ 30 phút, kim phút chỉ vào số 6, kim giờ ở vị trí chính giữa số 2 và 3 ($2{,}5$).\\
Khoảng cách giữa hai kim: $6 - 2{,}5 = 3{,}5$ khoảng giờ.\\
Số đo góc tạo thành là:
$$3{,}5 \times 30^\circ = 105^\circ$$
}
\end{ex}

\begin{ex} %[4H2V07]
Cho hai góc kề nhau có tổng số đo bằng $90^\circ$. Biết góc thứ nhất lớn hơn góc thứ hai là $18^\circ$. Tính số đo của góc thứ nhất theo độ (chỉ ghi số).
\shortans{54}
\loigiai{
\textbf{Đáp số:} $54$\\
Số đo của góc thứ nhất là:
$$(90^\circ + 18^\circ) : 2 = 54^\circ$$
}
\end{ex}

\begin{ex} %[4H2V07]
Trong một ngày đêm (24 giờ), có bao nhiêu lần hai kim giờ và kim phút của đồng hồ tạo thành một góc vuông đúng vào các thời điểm đúng giờ (chẵn giờ)?
\shortans{4}
\loigiai{
\textbf{Đáp số:} $4$\\
Vào các thời điểm đúng giờ, hai kim tạo thành góc vuông khi kim phút chỉ số 12 và kim giờ chỉ số 3 hoặc số 9.\\
Trong 24 giờ có:
- 2 lần lúc 3 giờ (3 giờ sáng và 15 giờ chiều).
- 2 lần lúc 9 giờ (9 giờ sáng và 21 giờ tối).\\
Tổng cộng có $2 + 2 = 4$ lần.
}
\end{ex}
